% Imporditud: tools/import_eksperimendid.py
% Allikas: Eesti füüsikaolümpiaadi eksperimendi ülesannete
%   kogu (Rael Kalda, 2023), ülesanne Ü71, lahendus L71.
\setAuthor{Hans Daniel Kaimre}
\setRound{lõppvoor}
\setYear{2019}
\setNumber{P E2}
\setDifficulty{2}
\setTopic{dünaamika}
\setVahendid{klots, dünamomeeter, joonlaud}
\setPunktid{12}
\setAllikas{Eksperimendi kogu Ü71, lahendus lk 65}
\setLahendusseis{toores}

\prob{Klots}
Leida klotsi teibiga kaetud tahu ning laua vaheline hõõrdetegur.
Raskuskiirendus $g = \SI{9,8}{m/s^2}$.

\textit{Märkus:} Hõõrdetegur $\mu = \dfrac{F_h}{F_r}$, kus $F_h$ on hõõrdejõud
ning $F_r$ rõhumisjõud.

\hint

\solu
Esmalt leiame klotsi massi. Asetame selle lauale pikemale küljele. Konksust dü-

namomeetriga õige pisut klotsi kergitades saame kangi, mille jaoks kehtib jõumo-

mentide tasakaal mg b = F1b, kus $F_1$ on dünamomeetri näit ning b klotsi külje pikkus. Klotsi mass avaldub kui m = 2F1/g.

Nüüd tõmbame dünamomeetriga klotsi mööda laua pinda ühtlase kiirusega. Tõm-

bejõud (dünamomeetri näit $F_2$) ja hõõrdejõud on tasakaalus, s.t. $F_2$ = mg$\mu$. Tehes

siin asenduse m = 2F1/g, avaldame hõõrdeteguri:

$\mu$ = $F_2$ 2F1

Asendades mõõtmisel saadud tulemused võrrandisse, saame et $\mu \approx$ ...
\probend
